\section{Introduction}
\label{sec:introduction}

Open-source package registries such as PyPI are a primary vector for software
supply-chain attacks: adversaries publish packages that masquerade as popular
libraries---through typosquatting, name squatting, or dependency
confusion---and execute malicious code during installation.
Studies of real-world registries have documented thousands of such
incidents~\cite{duan2021supply, ohm2020backstabber}, and the attack surface
keeps growing because anyone can publish, and installation scripts
(\texttt{setup.py}) are arbitrary code. Manual review does not scale, which
has motivated static-analysis heuristics and, more recently, large language
models (LLMs) that classify source code as malicious or benign.

LLMs are promising but imperfect reviewers: an 8B-scale instruct model is
strong at recognizing overt malware, yet it misses subtle cases such as
obfuscated droppers or typosquats whose installation script looks, on its
face, like any other package~\cite{ease2025rag}. One prominent remedy is
retrieval-augmented generation (RAG): before deciding, show the model
similar examples retrieved from a corpus of previously seen
malware~\cite{ease2025rag, crag2024}. Recent work applying RAG to malicious
PyPI package detection, most directly the EASE 2025 study ``Detecting
Malicious Source Code in PyPI Packages with LLMs: Does RAG Come in
Handy?''~\cite{ease2025rag}, reports that retrieved context improves detection,
but that study coupled retrieval to particular knowledge sources (external
threat-intelligence collections such as YARA rules and GitHub advisories) and
a single retrieval mechanism, leaving an important question open: \emph{how
much of the benefit comes from retrieval itself, and how much from how
retrieval is done?}

\begin{figure}[!ht]
    \centering
    \fbox{%
    \begin{minipage}[c]{0.995\linewidth}
        \centering
        \emph{Research question.}
        Given a fixed code-only knowledge base and a fixed LLM,
        which retrieval strategy---none, dense, sparse, hybrid,
        behavior-based, or contrastive---best detects malicious PyPI packages,
        and how do the strategies trade off recall, false alarms,
        and explanation faithfulness?
    \end{minipage}%
    }
    \caption{The research question addressed in this paper.}
    \label{fig:rq}
    \Description{Boxed statement of the paper's research question comparing
    retrieval strategies for malicious package detection.}
\end{figure}

This paper reports a controlled comparison of seven pipelines that differ
only in their retrieval strategy, evaluated on the same dataset used by the
EASE 2025 study, with a single fixed LLM (\texttt{Llama-3.1-8B-Instruct})
and a knowledge base built purely from the dataset's own training partition.
Our results show that:

\begin{itemize}
    \item \textbf{Any code-level retrieval is a large win.} Dense or sparse
    retrieval lifts malicious-class F1 from 0.888 (no RAG) to 0.968--0.972
    (+8 points), driven almost entirely by a 17-point recall gain
    (0.798~$\rightarrow$~0.972).
    \item \textbf{Sparse retrieval suffices.} Plain BM25 over source tokens
    matches a 4B embedding model within 0.4 F1 points, and hybrid fusion of
    the two does not improve on the better single retriever.
    \item \textbf{Contrastive retrieval minimizes false alarms.} Retrieving
    benign counterexamples alongside malicious ones lowers the false-positive
    rate from 0.8--1.2\% to 0.1\% while preserving 0.948 recall and
    0.996 precision.
    \item \textbf{AST behavior summaries are lossy.} Representing code as
    extracted behavior summaries for retrieval \emph{hurts}
    (F1 0.878), below even the no-RAG baseline.
\end{itemize}

\section{Motivation}
\label{sec:motivation}

Two practical concerns motivate the comparison. First, deployment cost:
if a lexicon index suffices, teams can skip embedding infrastructure,
vector stores, and the associated hosting budget. Second, triage cost:
security-reviewing LLM-predicted alarms is expensive, so the false-positive
rate per millisecond-of-recall may matter more than raw F1. By measuring
recall, false alarms, coverage (parse failures), and explanation grounding
side by side, we aim to give practitioners---and follow-up researchers---a
basis to choose a retrieval strategy on purpose rather than by habit.

\section{Contributions}
\label{sec:contributions}

The main contributions of this paper are:
\begin{itemize}
    \item A controlled, reproducible comparison of six RAG variants plus a
    zero-shot baseline (E0--E6) for LLM-based detection of malicious PyPI
    packages on the EASE 2025 dataset, isolating retrieval strategy as the
    only varying factor.
    \item Evidence that sparse retrieval matches dense retrieval on this
    task (0.968 vs.\ 0.972 F1) with no learned components.
    \item Evidence that contrastive dual-pool retrieval reduces false alarms
    by an order of magnitude at modest recall cost.
    \item A negative result for behavior-summary retrieval, quantifying what
    is lost when code is reduced to AST-derived descriptions.
    \item An analysis of replication-grade failure modes: samples no method
    detects (metadata-cloned typosquats), recurring false positives, and
    structured-output parse failures, with per-sample raw outputs and a
    reproducible evaluation pipeline released as a replication package.
\end{itemize}

\section{Paper Structure}
\label{sec:structure}

The remainder of this paper is structured as follows.
Section~\ref{sec:related-work} reviews related work on supply-chain attacks,
LLMs for malicious-code detection, and retrieval-augmented methods.
Section~\ref{sec:methodology} details the dataset, the seven retrieval
strategies, the LLM configuration, and the evaluation metrics.
Section~\ref{sec:results} presents the main results, fairness-oriented
faithfulness statistics, and an error analysis. Section~\ref{sec:discussion}
interprets the findings, discusses limitations and threats to validity, and
outlines future work. Section~\ref{sec:conclusion} concludes.
\endinput
